\documentclass{article}
\usepackage{xepersian}
\settextfont{Arial}
\usepackage{float}
\usepackage{graphicx}
\begin{document}
\pagenumbering{gobble}
\tableofcontents    
\newpage
\pagenumbering{arabic}
\section*{مقدمه}
\addcontentsline{toc}{section}{مقدمه}
\begin{flushright}
پس از بررسی پروژه را به چهار بخش تقسیم کردیم که به صورت زیر است:
	\vspace{2mm}
	
	\textbf{بخش اول:} پیاده‌سازی الگوریتم انتگرال‌گیری عددی سیمپسون مرکب با ساختار برداری،ارزیابی رفتار و همگرایی تابع.
	
	\textbf{بخش دوم:} پیاده‌سازی انتگرال‌گیری عددی تودرتو برای محاسبه‌ی میانگین تابع، و حل معادله‌ی F(x) جهت یافتن نقطه‌ی $c1$ برای با استفاده از روش نیوتن-رافسون و تحلیل حساسیت به حدس اولیه.
	
	\textbf{بخش سوم:}  بازسازی تابع با استفاده از روش‌های درون‌یابی، مشتق‌گیری عددی از تابع تقریب ‌زده‌شده، و محاسبه‌ی نقطه‌ی  $c2 $ با به ‌کارگیری روش دوبخشی  .
	
 \textbf{بخش چهارم: }استخراج پاسخ‌های دقیق با میپل ، قضاوت عددی و تحلیل خطای مطلق و نسبی الگوریتم‌ها.
 
 
\section{بخش اول}
\subsection{روش عددی انتخاب ‌شده}
برای تابع مورد نظر و انتگرال گرفتن از آن سه روش مستطیلی، ذوزنقه‌ای و سیمپسون را داشتیم . در روش ذوزنقه‌ای و مستطیلی تابع را با خط راست بین نقاط تقریب می‌زنیم که خطای همگرایی‌شان کند است. پس از روش سیمپسون استفاده می‌کنیم که بین هر سه نقطه‌ی متوالی یک سهمی درجه ۲ رد  می‌کند بنابراین برای تابع ما یعنی $ \cos(e^t) $ که مشتق آن پیوسته است خطا با نرخ $O(h^4)$ کاهش می‌یابد یعنی با هر بار دوبرابر کردن تعداد زیربازه‌ها خطا حدود ۱۶ برابر کوچک‌تر می‌شود در حالی‌که در روش ذوزنقه‌ای فقط ۴ برابر کوچک می‌شود پس برای رسیدن به دقت یکسان، در روش سیمپسون تعداد محاسبات کمتری داریم و از آنجا که $ \cos(e^t) $ در بازه مورد نظر هیچ نوسان شدیدی ندارد، پس شرایط ایدئال برای استفاده از روش سیمپسون فراهم است.
\subsection{ساختار کدنویسی و پیاده‌سازی در متلب}
\begin{enumerate}
	\item \textbf{تابع \texttt{simpsonf}:} تابع پایه انتگرال‌گیری عددی که ورودی برداری را می‌پذیرد و وزن‌های روش سیمپسون ($1, 4, 2, 4, ..., 1$) را اعمال می‌کند.
	\item \textbf{تابع \texttt{computeF}:} تابع اصلی دریافت‌کننده $x$ و تعداد زیربازه‌ها ($n$). این تابع حد بالای انتگرال را برابر $\tan(x)$ قرار داده و مقدار $F(x)$ را محاسبه می‌کند.
	\item \textbf{مقدار پیش‌فرض $n$:} تعداد زیربازه‌ها به‌صورت پیش‌فرض روی $n = 320$ تنظیم شد تا دقتی در حدود $10^{-12}$ برای محاسبات بعدی فراهم شود.
\end{enumerate}
\newpage
\subsection{بررسی همگرایی}
برای بررسی دقت روش پیاده‌سازی شده تابع $F(x)$ را در نقطه $x=0.25$ ارزیابی کردیم:
\begin{table}[H]
	\centering
	\begin{tabular}{|c|c|c|}
		\hline
		\textbf{تفاضل متوالی $|F_n - F_{n/2}|$}&\textbf{مقدار محاسباتی $F(0.25)$} &\textbf{تعداد زیربازه‌ها ($n$)}   \\ \hline
		— & $1.106425388974$ &$20$  \\ \hline
		 $1.475 \times 10^{-10}$& $1.106425388827$ & $40$ \\ \hline
		 $9.224 \times 10^{-12}$ & $1.106425388818$ &$80$ \\ \hline
		$5.764 \times 10^{-13}$& $1.106425388817$ &$160$   \\ \hline
		 $3.619 \times 10^{-14}$& $1.106425388817$ & $320$ \\ \hline
	\end{tabular}
	\caption{نتایج همگرایی $F(0.25)$}
	\label{tab:convergence_F025}
\end{table}
\paragraph{تحلیل خطای همگرایی:}همان‌طور که در جدول مشاهده می‌شود، با دو برابر شدن تعداد زیربازه‌ها ($n \to 2n$)، خطای برآوردی تقریباً $16$ برابر کاهش می‌یابد. این رفتار کاهش خطا دقیقا با تئوری خطای روش سیمپسون مرکب ($O(h^4)$) مطابقت کامل دارد و صحت پیاده‌سازی را تایید می‌کند.
\subsection{محاسبه نقاط نمونه و رفتار نموداری تابع}
برای ارزیابی رفتار کلی تابع در بازه داده شده مقادیر $F(x)$ روی نقاط نمونه با $ n=320$ استخراج شد
\begin{figure}[htbp]
	\centering
	\includegraphics[width=1.0\textwidth]{fig1.png}
	\label{fig:my_image}
\end{figure}
\newpage
\section{بخش دوم}
در این بخش، هدف محاسبات عددی مربوط به قضیه‌ی مقدار میانگین انتگرال برای تابع $F(x)$ در بازه‌ی $[a, b] = [0, 0.5]$ می‌باشد. طبق این قضیه نقطه‌ای مانند $c \in [a, b]$ وجود دارد به طوری که:
\begin{equation}
	\frac{1}{b-a} \int_a^b F(t) \, dt = F(c)
\end{equation}
 پیاده‌سازی عددی این رابطه را طی دو مرحله انجام دادیم. در مرحله اول ، برای محاسبه‌ی سمت چپ معادله از روش انتگرال گیری سیمپسون مرکب استفاده شده. از آنجا که خود تابع $F(t)$ شامل انتگرال گیری عددی است ، فرآیند محاسبه به صورت انتگرال‌گیری تو در تو انجام می‌شود و در مرحله دوم برای یافتن ریشه c در معادله $g(x) = F(x) - \bar{F} = 0$ از روش نیوتن-رافسون استفاده شد.
 \subsection{ساختار کدنویسی و پیاده‌سازی در متلب}
 \begin{enumerate}
 	\item \textbf{فراخوانی توابع پایه:} استفاده از تابع computef و simpsonf بخش اول.
 	\item \textbf{محاسبه‌ی انتگرال بیرونی:} مقدار میانگین انتگرالی $\bar{F}$ با اعمال روش سیمپسون روی بازه‌ی $[0, 0.5]$ با تعداد زیربازه‌های $n_{outer}$ محاسبه شد.
 	\item \textbf{تابع \texttt{newtonRaphson}:} الگوریتم نیوتن-رافسون که مشتق تابع را به روش تفاضل مرکزی با گام $h = 10^{-6}$ محاسبه کرده و تا رسیدن به تلرانس $10^{-10}$ تکرار می‌شود.
 \end{enumerate}
 \subsection{ ریشه‌یابی با نیوتن-رافسون}
  \begin{table}[H]
  	\centering
  	\begin{tabular}{ccc}
  		\hline
  		\textbf{$g(x)$}& \textbf{$x$} & \textbf{تکرار}  \\
  		\hline
  		$-1.323 \times 10^{-3}$ & \lr{0.203166786529} & $1$ \\
  		 $-7.996 \times 10^{-6}$& \lr{0.206950242268} & $2$ \\
  		 $-3.012 \times 10^{-10}$ & \lr{0.206973393074} & $3$\\
  		$0.000 \times 10^{0}$ & \lr{0.206973393946} & $4$ \\
  		\hline
  	\end{tabular}
  	  	\caption{تکرارهای روش نیوتن-رافسون برای یافتن $c$ (حدس اولیه $x_0 = 0.25$)}
  	\label{tab:newton-iterations}
  \end{table}
  میانگین انتگرالی روی $[0, 0.5]$ (با $n_{outer}=40$) برابر است با: $1.092644024396$ 
  
  و ریشه‌ی یافت‌شده $c = {0.206973393946}$ و $F(c) = 1.092644024396$ که دقیقا برابر با مقدار میانگین محاسبه‌شده است.
 \subsection{بررسی همگرایی و حساسیت به حدس اولیه}
 \begin{table}[H]
 	\centering

 	\begin{tabular}{|c|c|c|}
 		\hline
 		\textbf{نقطه‌ی ریشه ($c$)} & \textbf{مقدار میانگین ($\bar{F}$)} & \textbf{تعداد زیربازه‌ها ($n_{outer}$)} \\ \hline
 		$0.206970879871$ & $1.092643156089$ & $10$ \\ \hline
 		$0.206973245548$ & $1.092643973143$ & $20$ \\ \hline
 		$0.206973393946$ & $1.092644024396 $ & $40$ \\ \hline
 		$0.206973403229$ & $1.092644027602$ & $80$ \\ \hline
 		$0.206973403810$ & $1.092644027803 $ & $160$ \\ \hline
 	\end{tabular}
 	 \caption{نتایج همگرایی $\bar{F}$ و نقطه‌ی $c$ نسبت به $n_{outer}$}
 	\label{tab:convergence_person2}
 \end{table}
 \paragraph{بررسی حساسیت به حدس اولیه و شکست همگرایی نزدیک مرز بازه:}
 برای اکثر حدس‌های اولیه ($0.01 , 0.05,0.1,0.25,0.40$) نیوتن-رافسون درست به$0.206973393946$ همگرا می‌شود اما در نزدیکی مرز بازه یعنی نقطه $0.49$ نیوتن-رافسون به یک ریشه‌ی کاملاً غلط و خارج از بازه همگرا می‌شود زیرا نزدیک $ x=0.5$ شیب$ tan(x)$ به شدت تند شده و مشتق عددی g در آن ناحیه رفتار رفتار حساسی داشته و قدم اول نیوتن-رافسون می‌تواند خیلی بزرگ باشد و از بازه‌ی مجاز بیرون بپرد؛ وقتی یک بار از بازه خارج شد، دیگر به سمت ریشه‌ی درست برنمی‌گردد و به یک ریشه‌ی نامربوط همگرا می‌شود.
 \subsection{رفتار نموداری}
 \begin{figure}[htbp]
 	\centering
 	\includegraphics[width=1.0\textwidth]{fig2.png}
 	\label{fig:my_image}
 \end{figure}
 \newpage
 \section{بخش سوم}
 
\end{flushright}
\end{document}